\documentclass[11pt]{letter}

\usepackage[margin=0.5in]{geometry}
\usepackage{hyperref}

\begin{document}

\begin{letter}{}

\opening{Dear Hiring Team,}

I am writing to apply for the Junior Software Engineer (Remote) position at PolicyMe. With an M.Sc. in Computer Science from the University of Calgary and hands-on experience building Python software, Django REST APIs, data pipelines, and AI-enabled automation workflows, I am excited by the opportunity to contribute to a remote-first team building practical AI features for insurance.

What interests me most about this role is the mix of software engineering, product ownership, and applied AI. PolicyMe's posting describes work across backend services, AI model and API integration, data-driven product capabilities, code review, architecture discussions, and end-to-end feature development. That combination fits the kind of engineering I enjoy most: building useful systems, understanding the data and users behind them, and improving reliability through careful debugging, documentation, and iteration.

My strongest language is Python, and I have used it across backend development, scientific software, machine-learning workflows, and automation projects. As a Back-End Developer at Asr Gooyesh Pardaz, I built Django REST APIs backed by PostgreSQL for NLP, text-to-speech, and speech-to-text services. I implemented authentication and verification flows, payment-related backend logic, and database-backed service-plan features. That experience gave me practical grounding in API design, relational data modeling, integration work, and writing maintainable backend code.

During my graduate research with BigGeo, Mitacs, and the University of Calgary, I built Python and C++ software for processing, storing, and retrieving large Earth observation datasets. The work required designing data structures, automating preprocessing pipelines, debugging incorrect outputs, measuring performance, and improving one C++ encoding workload from 233 seconds to 26 seconds through multithreading. Although the domain was geospatial research, the engineering challenges were very transferable to product software: making complex data usable, keeping workflows reproducible, and turning ambiguous requirements into reliable systems.

I have also been building an agentic job-search automation platform using Docker Compose, n8n, LLM-based workflows, Google Sheets and Drive integrations, SearXNG, and Telegram. This personal project is especially relevant to PolicyMe's AI-focused product work because it combines workflow automation, practical LLM use, external integrations, structured data, and user-facing decision support. It has strengthened my interest in AI systems that save people time in real workflows rather than existing only as demos.

I would bring strong Python skills, backend/API experience, PostgreSQL familiarity, AI/ML foundations, and a careful, product-minded engineering style. React, Redux, and LangChain/LangGraph are areas where I would still be growing, but I am comfortable learning new tools quickly and contributing first where my Python, backend, data, and automation experience are strongest. I would be glad to help PolicyMe build intelligent, reliable software while continuing to grow across the full stack.

Thank you for considering my application. I would welcome the opportunity to discuss how my background in backend engineering, AI workflows, and data-intensive software can contribute to PolicyMe's engineering team.

\closing{Sincerely,

Amir Mirzai Golpayegani

\href{mailto:amirgolpaa24@gmail.com}{amirgolpaa24@gmail.com}}

\end{letter}

\end{document}
